% =====================================================================
%  SEÇÃO DE MÉTODO — Revisão Sistemática da Literatura (RSL)
%  Inserida em isys_rsl.tex via  % =====================================================================
%  SEÇÃO DE MÉTODO — Revisão Sistemática da Literatura (RSL)
%  Inserida em isys_rsl.tex via  % =====================================================================
%  SEÇÃO DE MÉTODO — Revisão Sistemática da Literatura (RSL)
%  Inserida em isys_rsl.tex via  % =====================================================================
%  SEÇÃO DE MÉTODO — Revisão Sistemática da Literatura (RSL)
%  Inserida em isys_rsl.tex via  \input{metodo_rsl}
%
%  CLASSE: sbc2025 (apalike-sol). Sem comandos ABNT (\quadro, \citeonline).
%  PACOTES: usa apenas tabular/\hline padrão — não exige booktabs.
%
%  PARÂMETROS A CONFIRMAR (procure por "CONFIRMAR" abaixo):
%   - Janela temporal (ano inicial)
%   - String final após calibração-piloto
%   - Limiares de qualidade
%
%  REFERÊNCIAS .bib:
%   NOVAS (adicionadas a referencias.bib): kitchenham2007, petersen2015,
%         page2021, wohlin2014.
%   NOTA: Kitchenham2004 (mapeamento) já existe no .bib, mas é obra distinta;
%         o Método cita Kitchenham & Charters (2007) = kitchenham2007.
%
%  OBS.: se a classe for de duas colunas e alguma tabela estourar,
%        troque \begin{table} por \begin{table*}.
% =====================================================================

\section{Método}
\label{sec:metodo}

Este trabalho caracteriza-se como uma Revisão Sistemática da Literatura (RSL), método secundário cujo propósito é identificar, avaliar e sintetizar as evidências científicas disponíveis sobre um dado fenômeno de forma reprozível. O protocolo foi planejado e conduzido segundo as diretrizes de Kitchenham e Charters \citeyearpar{kitchenham2007}, complementadas pelas orientações para mapeamento sistemático de Petersen et al. \citeyearpar{petersen2015}. O relato da pesquisa obedece às recomendações do framework PRISMA 2020 \citep{page2021}.

Como as questões levantadas são de caracterização, e não explicativas, optou-se por um desenho \textbf{híbrido}. Isso significa que a busca e a seleção de literatura seguem o rigor procedimental de uma RSL, mas a síntese dos resultados assume a abrangência de um mapeamento de campo. A execução organiza-se em três etapas: planejamento (definição do protocolo e critérios), condução (busca, seleção e extração) e relato (síntese e discussão). Visando proteger a reprodutibilidade do estudo e refrear o viés do pesquisador, todo o protocolo foi consolidado \textit{a priori}.

\subsection{Objetivo e Questões de Pesquisa}
\label{subsec:objetivo_qp}

O objetivo desta revisão é identificar como a Inteligência Artificial Generativa e os Modelos de Linguagem de Grande Porte (\textit{Large Language Models} -- LLMs) vêm sendo empregados para destravar o acesso a dados governamentais abertos e impulsionar a transparência pública. O esforço ancora-se na área de Sistemas de Informação, operando exatamente na interseção entre governo digital, política de dados abertos (\textit{Open Government Data} -- OGD) e apoio à decisão. A partir dessa diretriz central, derivam-se as cinco questões de pesquisa detalhadas na Tabela~\ref{tab:rqs}.

\begin{table*}[t]
\centering
\caption{Questões de pesquisa e foco de extração}
\label{tab:rqs}
\small
\begin{tabular}{|c|p{4.6cm}|p{4.0cm}|}
\hline
\textbf{ID} & \textbf{Questão de pesquisa} & \textbf{Foco da extração} \\ \hline
RQ1 & Quais aplicações de IA Generativa são utilizadas no acesso a dados governamentais abertos e na transparência pública? & Tipo de aplicação / caso de uso \\ \hline
RQ2 & Quais tipos de dados governamentais são analisados? & Natureza e domínio dos dados \\ \hline
RQ3 & Quais LLMs são utilizados? & Modelo(s), porte, aberto/proprietário \\ \hline
RQ4 & Quais benefícios são reportados? & Ganhos de acesso e transparência \\ \hline
RQ5 & Quais desafios e limitações são reportados? & Barreiras técnicas, éticas e de adoção \\ \hline
\end{tabular}
\end{table*}

\subsection{Estratégia de Busca}
\label{subsec:busca}

A varredura automática ocorreu em quatro bases de amplo reconhecimento no domínio da Computação: Scopus, Web of Science, ACM Digital Library e IEEE Xplore. Essa seleção equilibra indexadores multidisciplinares (Scopus, WoS) e repositórios de nicho tecnológico (ACM, IEEE), mitigando o risco de omissão de publicações chave.

Os termos da consulta foram arquitetados a partir das cinco dimensões da estrutura PICOC (Tabela~\ref{tab:picoc}), garantindo que a \textit{string} não recaísse isoladamente sobre a Intervenção tecnológica. O eixo de Intervenção mapeia o maquinário propriamente dito: IA Generativa, modelos de fundação, \textit{transformers}, arquiteturas RAG, agentes conversacionais e LLMs comerciais ou abertos. Em contrapartida, os blocos de População, Desfecho e Contexto amarram o escopo da ferramenta ao setor público, filtrando estritamente estudos sobre portais governamentais, dados legislativos e políticas de transparência digital.

A \textit{string} de referência consta na Figura~\ref{fig:string} e foi adaptada às regras sintáticas de cada base no momento do disparo. Para assegurar aderência ao critério de idioma (CI4), aplicou-se uma variação do comando em português à Biblioteca Digital da SBC/SOL e demais plataformas lusófonas. Com o intuito de distinguir a real escassez de produção acadêmica de meros erros de digitação na busca, a estrutura sintática exata, a data e o volume retornado por base foram documentados no repositório do projeto e recuperados na Seção~\ref{sec:resultados}.

\begin{table*}[t]
\centering
\caption{Termos da busca segundo a estrutura PICOC}
\label{tab:picoc}
\small
\begin{tabular}{|p{2.4cm}|p{6.0cm}|}
\hline
\textbf{Elemento} & \textbf{Termos} \\ \hline
População & Setor público; portais e dados governamentais \\ \hline
Intervenção & IA Generativa; LLMs; modelos de linguagem \\ \hline
Comparação & Não aplicável (revisão de caracterização) \\ \hline
Desfecho & Acesso à informação; transparência; benefícios e desafios \\ \hline
Contexto & Governo digital; \textit{e-government}; OGD \\ \hline
\end{tabular}
\end{table*}

\begin{figure*}[t]
\centering
\begin{minipage}{0.95\textwidth}
\begin{small}
\begin{verbatim}
("generative artificial intelligence" OR "generative AI"
 OR "large language model*" OR "LLM*" OR "foundation model*"
 OR transformer* OR GPT* OR ChatGPT OR Llama OR Mistral OR Gemini
 OR "retrieval-augmented generation" OR RAG
 OR chatbot* OR "conversational agent*")
AND
("open government data" OR "open data" OR "government data"
 OR "public sector data" OR "legislative data" OR "parliamentary data"
 OR "freedom of information" OR FOIA OR "access to information"
 OR transparency OR accountability
 OR "digital government" OR "e-government" OR "open government"
 OR "data portal*")
\end{verbatim}
\end{small}
\end{minipage}
\caption{String de busca de referência, com dois blocos conectados por \texttt{AND}: o primeiro reúne termos de IA Generativa/LLMs (Intervenção); o segundo, termos de dados governamentais abertos, governo digital e transparência (População, Desfecho, Contexto). A sintaxe é adaptada a cada base no momento da execução.}
\label{fig:string}
\end{figure*}

Antes da execução definitiva, a \textit{string} passou por calibração-piloto. Ela foi confrontada contra um acervo de controle previamente mapeado pela literatura: caso algum desses artigos essenciais não fosse detectado, o escopo de palavras-chave sofreria expansão imediata. % CONFIRMAR: liste 3-5 artigos de controle usados na calibração.

O intervalo temporal abrange de janeiro de 2019 até a data da pesquisa. % CONFIRMAR ano inicial.
A justificativa para o corte repousa no fato de a consolidação dos grandes modelos pré-treinados ter ocorrido sistematicamente a partir de 2019. O esforço apoiou-se exclusivamente no retorno das bases de dados; técnicas auxiliares de prospecção referencial em cascata (\textit{snowballing} para frente e para trás) \citep{wohlin2014} não foram acionadas neste estágio, configurando uma limitação metodológica discutida com maior profundidade na Seção~\ref{sec:ameacas}.

\subsection{Critérios de Seleção e Processo de Triagem}
\label{subsec:selecao}

O filtro do \textit{corpus} bibliográfico baseou-se nos parâmetros expostos na Tabela~\ref{tab:criterios}. O protocolo é excludente: a aprovação final de um artigo exige a conformidade simultânea a todos os critérios de inclusão e a imunidade absoluta a qualquer um dos critérios de exclusão.

\begin{table}[ht]
\centering
\caption{Critérios de inclusão e exclusão}
\label{tab:criterios}
\small
\begin{tabular}{|c|p{6.4cm}|}
\hline
\textbf{ID} & \textbf{Critério} \\ \hline
CI1 & Aplica ou investiga IA Generativa/LLMs em contexto de dados abertos, transparência pública ou governo digital \\ \hline
CI2 & Estudo primário revisado por pares (conferência ou periódico) \\ \hline
CI3 & Publicado na janela temporal definida \\ \hline
CI4 & Texto em inglês ou português \\ \hline
CI5 & Texto completo acessível \\ \hline
CE1 & Aborda apenas um dos eixos (somente IA \textit{ou} somente governo aberto), sem a interseção \\ \hline
CE2 & Estudo secundário (outra revisão, mapeamento ou \textit{survey}) -- considerado apenas como referência para trabalhos relacionados \\ \hline
CE3 & Literatura cinza, editorial, resumo estendido, pôster, \textit{keynote} ou tutorial \\ \hline
CE4 & Duplicata \\ \hline
CE5 & Versão preliminar de estudo posteriormente estendido (mantém-se a versão mais completa) \\ \hline
\end{tabular}
\end{table}

O funil de seleção operou em duas etapas. A triagem inicial ocorreu por meio da leitura rápida de títulos e resumos. Em seguida, os trabalhos pré-selecionados avançaram para a avaliação na íntegra. Todo o referencial foi exportado para o Zotero (responsável pela deduplicação primária) e posteriormente importado para a plataforma Rayyan, onde ocorreu a triagem manual.

Uma única pesquisadora conduziu a revisão. Para mitigar o viés interpretativo individual, todos os parâmetros de inclusão e exclusão foram engessados formalmente no projeto de pesquisa. Além disso, as exclusões no Rayyan exigiram a marcação explícita do rótulo motivador correspondente ao critério CE, gerando uma trilha de auditoria rastreável. O diagrama PRISMA, ilustrado na seção de resultados, consolida todo o fluxo das exclusões.

\subsection{Avaliação de Qualidade}
\label{subsec:qualidade}

Os manuscritos sobreviventes à leitura integral foram então submetidos a um crivo de qualidade focado no rigor científico das evidências. Cada métrica da Tabela~\ref{tab:qualidade} comporta três valores possíveis: $0$ (não atende), $0{,}5$ (atende parcialmente) e $1$ (atende plenamente), culminando em uma pontuação máxima de 5.

Estabeleceu-se o limiar de $2{,}5$ como nota de corte analítica. Trabalhos que igualam ou superam essa barreira compõem o núcleo da síntese informacional sem ressalvas. Aqueles com desempenho inferior não sofrem descarte sumário, mas são apartados em um bloco interpretativo paralelo, exigindo ceticismo deliberado na fase de Discussão.

Esse recurso avaliativo opera também como mecanismo compensatório para publicações em veículos de menor tradição editorial (como simpósios de iniciação ou anais de triagem branda). Pelo critério CI2, estudos primários em eventos oficiais não são descartados na fase inicial. É a matriz de qualidade (notadamente os fatores AQ4 e AQ5) que penaliza a fragilidade argumentativa dessas produções, desidratando sua influência na conclusão final sem violar as regras da filtragem sistemática.

\begin{table}[ht]
\centering
\caption{Critérios de avaliação de qualidade}
\label{tab:qualidade}
\small
\begin{tabular}{|c|p{6.4cm}|}
\hline
\textbf{ID} & \textbf{Critério} \\ \hline
AQ1 & Os objetivos do estudo estão claramente declarados? \\ \hline
AQ2 & O contexto (domínio governamental e dados) está adequadamente descrito? \\ \hline
AQ3 & A abordagem de IA Generativa/LLM está descrita com detalhe suficiente para ser compreendida e replicada? \\ \hline
AQ4 & Os resultados e benefícios são reportados e sustentados por evidências? \\ \hline
AQ5 & As limitações e ameaças à validade são discutidas? \\ \hline
\end{tabular}
\end{table}

\subsection{Extração de Dados}
\label{subsec:extracao}

A captação de evidências seguiu um formulário rígido cujas entradas convergem diretamente para as cinco questões de pesquisa (Tabela~\ref{tab:extracao}). Como complemento à estruturação técnica, levantou-se o inventário bibliométrico padrão (autoria, ano, país e tipologia do veículo), bem como as matrizes metodológicas de validação presentes em cada artigo empírico. Todo o registro das abstrações foi conduzido pela revisora primária, mantendo a indexação rigorosa da informação recolhida com o arquivo de origem para evitar distorções de interpretação fora de contexto.

\begin{table*}[t]
\centering
\caption{Formulário de extração e vínculo com as questões de pesquisa}
\label{tab:extracao}
\small
\renewcommand{\arraystretch}{1.15}
\begin{tabular}{|p{3.6cm}|@{\hspace{0.6em}}p{6.0cm}@{\hspace{0.6em}}|c|}
\hline
\textbf{Campo} & \textbf{Descrição} & \textbf{RQ} \\ \hline
Aplicação / caso de uso & Como a IA Generativa é empregada & RQ1 \\ \hline
Dados governamentais & Tipo e domínio dos dados analisados & RQ2 \\ \hline
Modelo(s) & LLM/IA Generativa, porte, licença & RQ3 \\ \hline
Benefícios & Ganhos de acesso/transparência relatados & RQ4 \\ \hline
Desafios & Limitações técnicas, éticas e de adoção & RQ5 \\ \hline
Avaliação & Método de validação do estudo & --- \\ \hline
\end{tabular}
\end{table*}

\subsection{Síntese dos Dados}
\label{subsec:sintese}

O fechamento da revisão opera sob uma matriz dual. O primeiro pilar baseia-se na \textit{síntese quantitativa descritiva}, responsável por traduzir o volume documental em distribuições macro (como a linha do tempo de publicações, os LLMs mais citados em RQ3 e o espectro de dados analisados em RQ2). O segundo pilar é a \textit{síntese temática}, uma abordagem estritamente qualitativa. Nela, o esforço debruça-se sobre a codificação do texto bruto, extraindo e categorizando de forma orgânica as rotinas governamentais criadas (RQ1), os impactos sociais alcançados (RQ4) e as pedras de tropeço da adoção (RQ5). É o atrito analítico provocado pelo cruzamento dessas duas sínteses que permitirá, no momento da Discussão, organizar o caos de aplicações cívicas de Inteligência Artificial em uma taxonomia coesa e útil para a comunidade científica.

%
%  CLASSE: sbc2025 (apalike-sol). Sem comandos ABNT (\quadro, \citeonline).
%  PACOTES: usa apenas tabular/\hline padrão — não exige booktabs.
%
%  PARÂMETROS A CONFIRMAR (procure por "CONFIRMAR" abaixo):
%   - Janela temporal (ano inicial)
%   - String final após calibração-piloto
%   - Limiares de qualidade
%
%  REFERÊNCIAS .bib:
%   NOVAS (adicionadas a referencias.bib): kitchenham2007, petersen2015,
%         page2021, wohlin2014.
%   NOTA: Kitchenham2004 (mapeamento) já existe no .bib, mas é obra distinta;
%         o Método cita Kitchenham & Charters (2007) = kitchenham2007.
%
%  OBS.: se a classe for de duas colunas e alguma tabela estourar,
%        troque \begin{table} por \begin{table*}.
% =====================================================================

\section{Método}
\label{sec:metodo}

Este trabalho caracteriza-se como uma Revisão Sistemática da Literatura (RSL), método secundário cujo propósito é identificar, avaliar e sintetizar as evidências científicas disponíveis sobre um dado fenômeno de forma reprozível. O protocolo foi planejado e conduzido segundo as diretrizes de Kitchenham e Charters \citeyearpar{kitchenham2007}, complementadas pelas orientações para mapeamento sistemático de Petersen et al. \citeyearpar{petersen2015}. O relato da pesquisa obedece às recomendações do framework PRISMA 2020 \citep{page2021}.

Como as questões levantadas são de caracterização, e não explicativas, optou-se por um desenho \textbf{híbrido}. Isso significa que a busca e a seleção de literatura seguem o rigor procedimental de uma RSL, mas a síntese dos resultados assume a abrangência de um mapeamento de campo. A execução organiza-se em três etapas: planejamento (definição do protocolo e critérios), condução (busca, seleção e extração) e relato (síntese e discussão). Visando proteger a reprodutibilidade do estudo e refrear o viés do pesquisador, todo o protocolo foi consolidado \textit{a priori}.

\subsection{Objetivo e Questões de Pesquisa}
\label{subsec:objetivo_qp}

O objetivo desta revisão é identificar como a Inteligência Artificial Generativa e os Modelos de Linguagem de Grande Porte (\textit{Large Language Models} -- LLMs) vêm sendo empregados para destravar o acesso a dados governamentais abertos e impulsionar a transparência pública. O esforço ancora-se na área de Sistemas de Informação, operando exatamente na interseção entre governo digital, política de dados abertos (\textit{Open Government Data} -- OGD) e apoio à decisão. A partir dessa diretriz central, derivam-se as cinco questões de pesquisa detalhadas na Tabela~\ref{tab:rqs}.

\begin{table*}[t]
\centering
\caption{Questões de pesquisa e foco de extração}
\label{tab:rqs}
\small
\begin{tabular}{|c|p{4.6cm}|p{4.0cm}|}
\hline
\textbf{ID} & \textbf{Questão de pesquisa} & \textbf{Foco da extração} \\ \hline
RQ1 & Quais aplicações de IA Generativa são utilizadas no acesso a dados governamentais abertos e na transparência pública? & Tipo de aplicação / caso de uso \\ \hline
RQ2 & Quais tipos de dados governamentais são analisados? & Natureza e domínio dos dados \\ \hline
RQ3 & Quais LLMs são utilizados? & Modelo(s), porte, aberto/proprietário \\ \hline
RQ4 & Quais benefícios são reportados? & Ganhos de acesso e transparência \\ \hline
RQ5 & Quais desafios e limitações são reportados? & Barreiras técnicas, éticas e de adoção \\ \hline
\end{tabular}
\end{table*}

\subsection{Estratégia de Busca}
\label{subsec:busca}

A varredura automática ocorreu em quatro bases de amplo reconhecimento no domínio da Computação: Scopus, Web of Science, ACM Digital Library e IEEE Xplore. Essa seleção equilibra indexadores multidisciplinares (Scopus, WoS) e repositórios de nicho tecnológico (ACM, IEEE), mitigando o risco de omissão de publicações chave.

Os termos da consulta foram arquitetados a partir das cinco dimensões da estrutura PICOC (Tabela~\ref{tab:picoc}), garantindo que a \textit{string} não recaísse isoladamente sobre a Intervenção tecnológica. O eixo de Intervenção mapeia o maquinário propriamente dito: IA Generativa, modelos de fundação, \textit{transformers}, arquiteturas RAG, agentes conversacionais e LLMs comerciais ou abertos. Em contrapartida, os blocos de População, Desfecho e Contexto amarram o escopo da ferramenta ao setor público, filtrando estritamente estudos sobre portais governamentais, dados legislativos e políticas de transparência digital.

A \textit{string} de referência consta na Figura~\ref{fig:string} e foi adaptada às regras sintáticas de cada base no momento do disparo. Para assegurar aderência ao critério de idioma (CI4), aplicou-se uma variação do comando em português à Biblioteca Digital da SBC/SOL e demais plataformas lusófonas. Com o intuito de distinguir a real escassez de produção acadêmica de meros erros de digitação na busca, a estrutura sintática exata, a data e o volume retornado por base foram documentados no repositório do projeto e recuperados na Seção~\ref{sec:resultados}.

\begin{table*}[t]
\centering
\caption{Termos da busca segundo a estrutura PICOC}
\label{tab:picoc}
\small
\begin{tabular}{|p{2.4cm}|p{6.0cm}|}
\hline
\textbf{Elemento} & \textbf{Termos} \\ \hline
População & Setor público; portais e dados governamentais \\ \hline
Intervenção & IA Generativa; LLMs; modelos de linguagem \\ \hline
Comparação & Não aplicável (revisão de caracterização) \\ \hline
Desfecho & Acesso à informação; transparência; benefícios e desafios \\ \hline
Contexto & Governo digital; \textit{e-government}; OGD \\ \hline
\end{tabular}
\end{table*}

\begin{figure*}[t]
\centering
\begin{minipage}{0.95\textwidth}
\begin{small}
\begin{verbatim}
("generative artificial intelligence" OR "generative AI"
 OR "large language model*" OR "LLM*" OR "foundation model*"
 OR transformer* OR GPT* OR ChatGPT OR Llama OR Mistral OR Gemini
 OR "retrieval-augmented generation" OR RAG
 OR chatbot* OR "conversational agent*")
AND
("open government data" OR "open data" OR "government data"
 OR "public sector data" OR "legislative data" OR "parliamentary data"
 OR "freedom of information" OR FOIA OR "access to information"
 OR transparency OR accountability
 OR "digital government" OR "e-government" OR "open government"
 OR "data portal*")
\end{verbatim}
\end{small}
\end{minipage}
\caption{String de busca de referência, com dois blocos conectados por \texttt{AND}: o primeiro reúne termos de IA Generativa/LLMs (Intervenção); o segundo, termos de dados governamentais abertos, governo digital e transparência (População, Desfecho, Contexto). A sintaxe é adaptada a cada base no momento da execução.}
\label{fig:string}
\end{figure*}

Antes da execução definitiva, a \textit{string} passou por calibração-piloto. Ela foi confrontada contra um acervo de controle previamente mapeado pela literatura: caso algum desses artigos essenciais não fosse detectado, o escopo de palavras-chave sofreria expansão imediata. % CONFIRMAR: liste 3-5 artigos de controle usados na calibração.

O intervalo temporal abrange de janeiro de 2019 até a data da pesquisa. % CONFIRMAR ano inicial.
A justificativa para o corte repousa no fato de a consolidação dos grandes modelos pré-treinados ter ocorrido sistematicamente a partir de 2019. O esforço apoiou-se exclusivamente no retorno das bases de dados; técnicas auxiliares de prospecção referencial em cascata (\textit{snowballing} para frente e para trás) \citep{wohlin2014} não foram acionadas neste estágio, configurando uma limitação metodológica discutida com maior profundidade na Seção~\ref{sec:ameacas}.

\subsection{Critérios de Seleção e Processo de Triagem}
\label{subsec:selecao}

O filtro do \textit{corpus} bibliográfico baseou-se nos parâmetros expostos na Tabela~\ref{tab:criterios}. O protocolo é excludente: a aprovação final de um artigo exige a conformidade simultânea a todos os critérios de inclusão e a imunidade absoluta a qualquer um dos critérios de exclusão.

\begin{table}[ht]
\centering
\caption{Critérios de inclusão e exclusão}
\label{tab:criterios}
\small
\begin{tabular}{|c|p{6.4cm}|}
\hline
\textbf{ID} & \textbf{Critério} \\ \hline
CI1 & Aplica ou investiga IA Generativa/LLMs em contexto de dados abertos, transparência pública ou governo digital \\ \hline
CI2 & Estudo primário revisado por pares (conferência ou periódico) \\ \hline
CI3 & Publicado na janela temporal definida \\ \hline
CI4 & Texto em inglês ou português \\ \hline
CI5 & Texto completo acessível \\ \hline
CE1 & Aborda apenas um dos eixos (somente IA \textit{ou} somente governo aberto), sem a interseção \\ \hline
CE2 & Estudo secundário (outra revisão, mapeamento ou \textit{survey}) -- considerado apenas como referência para trabalhos relacionados \\ \hline
CE3 & Literatura cinza, editorial, resumo estendido, pôster, \textit{keynote} ou tutorial \\ \hline
CE4 & Duplicata \\ \hline
CE5 & Versão preliminar de estudo posteriormente estendido (mantém-se a versão mais completa) \\ \hline
\end{tabular}
\end{table}

O funil de seleção operou em duas etapas. A triagem inicial ocorreu por meio da leitura rápida de títulos e resumos. Em seguida, os trabalhos pré-selecionados avançaram para a avaliação na íntegra. Todo o referencial foi exportado para o Zotero (responsável pela deduplicação primária) e posteriormente importado para a plataforma Rayyan, onde ocorreu a triagem manual.

Uma única pesquisadora conduziu a revisão. Para mitigar o viés interpretativo individual, todos os parâmetros de inclusão e exclusão foram engessados formalmente no projeto de pesquisa. Além disso, as exclusões no Rayyan exigiram a marcação explícita do rótulo motivador correspondente ao critério CE, gerando uma trilha de auditoria rastreável. O diagrama PRISMA, ilustrado na seção de resultados, consolida todo o fluxo das exclusões.

\subsection{Avaliação de Qualidade}
\label{subsec:qualidade}

Os manuscritos sobreviventes à leitura integral foram então submetidos a um crivo de qualidade focado no rigor científico das evidências. Cada métrica da Tabela~\ref{tab:qualidade} comporta três valores possíveis: $0$ (não atende), $0{,}5$ (atende parcialmente) e $1$ (atende plenamente), culminando em uma pontuação máxima de 5.

Estabeleceu-se o limiar de $2{,}5$ como nota de corte analítica. Trabalhos que igualam ou superam essa barreira compõem o núcleo da síntese informacional sem ressalvas. Aqueles com desempenho inferior não sofrem descarte sumário, mas são apartados em um bloco interpretativo paralelo, exigindo ceticismo deliberado na fase de Discussão.

Esse recurso avaliativo opera também como mecanismo compensatório para publicações em veículos de menor tradição editorial (como simpósios de iniciação ou anais de triagem branda). Pelo critério CI2, estudos primários em eventos oficiais não são descartados na fase inicial. É a matriz de qualidade (notadamente os fatores AQ4 e AQ5) que penaliza a fragilidade argumentativa dessas produções, desidratando sua influência na conclusão final sem violar as regras da filtragem sistemática.

\begin{table}[ht]
\centering
\caption{Critérios de avaliação de qualidade}
\label{tab:qualidade}
\small
\begin{tabular}{|c|p{6.4cm}|}
\hline
\textbf{ID} & \textbf{Critério} \\ \hline
AQ1 & Os objetivos do estudo estão claramente declarados? \\ \hline
AQ2 & O contexto (domínio governamental e dados) está adequadamente descrito? \\ \hline
AQ3 & A abordagem de IA Generativa/LLM está descrita com detalhe suficiente para ser compreendida e replicada? \\ \hline
AQ4 & Os resultados e benefícios são reportados e sustentados por evidências? \\ \hline
AQ5 & As limitações e ameaças à validade são discutidas? \\ \hline
\end{tabular}
\end{table}

\subsection{Extração de Dados}
\label{subsec:extracao}

A captação de evidências seguiu um formulário rígido cujas entradas convergem diretamente para as cinco questões de pesquisa (Tabela~\ref{tab:extracao}). Como complemento à estruturação técnica, levantou-se o inventário bibliométrico padrão (autoria, ano, país e tipologia do veículo), bem como as matrizes metodológicas de validação presentes em cada artigo empírico. Todo o registro das abstrações foi conduzido pela revisora primária, mantendo a indexação rigorosa da informação recolhida com o arquivo de origem para evitar distorções de interpretação fora de contexto.

\begin{table*}[t]
\centering
\caption{Formulário de extração e vínculo com as questões de pesquisa}
\label{tab:extracao}
\small
\renewcommand{\arraystretch}{1.15}
\begin{tabular}{|p{3.6cm}|@{\hspace{0.6em}}p{6.0cm}@{\hspace{0.6em}}|c|}
\hline
\textbf{Campo} & \textbf{Descrição} & \textbf{RQ} \\ \hline
Aplicação / caso de uso & Como a IA Generativa é empregada & RQ1 \\ \hline
Dados governamentais & Tipo e domínio dos dados analisados & RQ2 \\ \hline
Modelo(s) & LLM/IA Generativa, porte, licença & RQ3 \\ \hline
Benefícios & Ganhos de acesso/transparência relatados & RQ4 \\ \hline
Desafios & Limitações técnicas, éticas e de adoção & RQ5 \\ \hline
Avaliação & Método de validação do estudo & --- \\ \hline
\end{tabular}
\end{table*}

\subsection{Síntese dos Dados}
\label{subsec:sintese}

O fechamento da revisão opera sob uma matriz dual. O primeiro pilar baseia-se na \textit{síntese quantitativa descritiva}, responsável por traduzir o volume documental em distribuições macro (como a linha do tempo de publicações, os LLMs mais citados em RQ3 e o espectro de dados analisados em RQ2). O segundo pilar é a \textit{síntese temática}, uma abordagem estritamente qualitativa. Nela, o esforço debruça-se sobre a codificação do texto bruto, extraindo e categorizando de forma orgânica as rotinas governamentais criadas (RQ1), os impactos sociais alcançados (RQ4) e as pedras de tropeço da adoção (RQ5). É o atrito analítico provocado pelo cruzamento dessas duas sínteses que permitirá, no momento da Discussão, organizar o caos de aplicações cívicas de Inteligência Artificial em uma taxonomia coesa e útil para a comunidade científica.

%
%  CLASSE: sbc2025 (apalike-sol). Sem comandos ABNT (\quadro, \citeonline).
%  PACOTES: usa apenas tabular/\hline padrão — não exige booktabs.
%
%  PARÂMETROS A CONFIRMAR (procure por "CONFIRMAR" abaixo):
%   - Janela temporal (ano inicial)
%   - String final após calibração-piloto
%   - Limiares de qualidade
%
%  REFERÊNCIAS .bib:
%   NOVAS (adicionadas a referencias.bib): kitchenham2007, petersen2015,
%         page2021, wohlin2014.
%   NOTA: Kitchenham2004 (mapeamento) já existe no .bib, mas é obra distinta;
%         o Método cita Kitchenham & Charters (2007) = kitchenham2007.
%
%  OBS.: se a classe for de duas colunas e alguma tabela estourar,
%        troque \begin{table} por \begin{table*}.
% =====================================================================

\section{Método}
\label{sec:metodo}

Este trabalho caracteriza-se como uma Revisão Sistemática da Literatura (RSL), método secundário cujo propósito é identificar, avaliar e sintetizar as evidências científicas disponíveis sobre um dado fenômeno de forma reprozível. O protocolo foi planejado e conduzido segundo as diretrizes de Kitchenham e Charters \citeyearpar{kitchenham2007}, complementadas pelas orientações para mapeamento sistemático de Petersen et al. \citeyearpar{petersen2015}. O relato da pesquisa obedece às recomendações do framework PRISMA 2020 \citep{page2021}.

Como as questões levantadas são de caracterização, e não explicativas, optou-se por um desenho \textbf{híbrido}. Isso significa que a busca e a seleção de literatura seguem o rigor procedimental de uma RSL, mas a síntese dos resultados assume a abrangência de um mapeamento de campo. A execução organiza-se em três etapas: planejamento (definição do protocolo e critérios), condução (busca, seleção e extração) e relato (síntese e discussão). Visando proteger a reprodutibilidade do estudo e refrear o viés do pesquisador, todo o protocolo foi consolidado \textit{a priori}.

\subsection{Objetivo e Questões de Pesquisa}
\label{subsec:objetivo_qp}

O objetivo desta revisão é identificar como a Inteligência Artificial Generativa e os Modelos de Linguagem de Grande Porte (\textit{Large Language Models} -- LLMs) vêm sendo empregados para destravar o acesso a dados governamentais abertos e impulsionar a transparência pública. O esforço ancora-se na área de Sistemas de Informação, operando exatamente na interseção entre governo digital, política de dados abertos (\textit{Open Government Data} -- OGD) e apoio à decisão. A partir dessa diretriz central, derivam-se as cinco questões de pesquisa detalhadas na Tabela~\ref{tab:rqs}.

\begin{table*}[t]
\centering
\caption{Questões de pesquisa e foco de extração}
\label{tab:rqs}
\small
\begin{tabular}{|c|p{4.6cm}|p{4.0cm}|}
\hline
\textbf{ID} & \textbf{Questão de pesquisa} & \textbf{Foco da extração} \\ \hline
RQ1 & Quais aplicações de IA Generativa são utilizadas no acesso a dados governamentais abertos e na transparência pública? & Tipo de aplicação / caso de uso \\ \hline
RQ2 & Quais tipos de dados governamentais são analisados? & Natureza e domínio dos dados \\ \hline
RQ3 & Quais LLMs são utilizados? & Modelo(s), porte, aberto/proprietário \\ \hline
RQ4 & Quais benefícios são reportados? & Ganhos de acesso e transparência \\ \hline
RQ5 & Quais desafios e limitações são reportados? & Barreiras técnicas, éticas e de adoção \\ \hline
\end{tabular}
\end{table*}

\subsection{Estratégia de Busca}
\label{subsec:busca}

A varredura automática ocorreu em quatro bases de amplo reconhecimento no domínio da Computação: Scopus, Web of Science, ACM Digital Library e IEEE Xplore. Essa seleção equilibra indexadores multidisciplinares (Scopus, WoS) e repositórios de nicho tecnológico (ACM, IEEE), mitigando o risco de omissão de publicações chave.

Os termos da consulta foram arquitetados a partir das cinco dimensões da estrutura PICOC (Tabela~\ref{tab:picoc}), garantindo que a \textit{string} não recaísse isoladamente sobre a Intervenção tecnológica. O eixo de Intervenção mapeia o maquinário propriamente dito: IA Generativa, modelos de fundação, \textit{transformers}, arquiteturas RAG, agentes conversacionais e LLMs comerciais ou abertos. Em contrapartida, os blocos de População, Desfecho e Contexto amarram o escopo da ferramenta ao setor público, filtrando estritamente estudos sobre portais governamentais, dados legislativos e políticas de transparência digital.

A \textit{string} de referência consta na Figura~\ref{fig:string} e foi adaptada às regras sintáticas de cada base no momento do disparo. Para assegurar aderência ao critério de idioma (CI4), aplicou-se uma variação do comando em português à Biblioteca Digital da SBC/SOL e demais plataformas lusófonas. Com o intuito de distinguir a real escassez de produção acadêmica de meros erros de digitação na busca, a estrutura sintática exata, a data e o volume retornado por base foram documentados no repositório do projeto e recuperados na Seção~\ref{sec:resultados}.

\begin{table*}[t]
\centering
\caption{Termos da busca segundo a estrutura PICOC}
\label{tab:picoc}
\small
\begin{tabular}{|p{2.4cm}|p{6.0cm}|}
\hline
\textbf{Elemento} & \textbf{Termos} \\ \hline
População & Setor público; portais e dados governamentais \\ \hline
Intervenção & IA Generativa; LLMs; modelos de linguagem \\ \hline
Comparação & Não aplicável (revisão de caracterização) \\ \hline
Desfecho & Acesso à informação; transparência; benefícios e desafios \\ \hline
Contexto & Governo digital; \textit{e-government}; OGD \\ \hline
\end{tabular}
\end{table*}

\begin{figure*}[t]
\centering
\begin{minipage}{0.95\textwidth}
\begin{small}
\begin{verbatim}
("generative artificial intelligence" OR "generative AI"
 OR "large language model*" OR "LLM*" OR "foundation model*"
 OR transformer* OR GPT* OR ChatGPT OR Llama OR Mistral OR Gemini
 OR "retrieval-augmented generation" OR RAG
 OR chatbot* OR "conversational agent*")
AND
("open government data" OR "open data" OR "government data"
 OR "public sector data" OR "legislative data" OR "parliamentary data"
 OR "freedom of information" OR FOIA OR "access to information"
 OR transparency OR accountability
 OR "digital government" OR "e-government" OR "open government"
 OR "data portal*")
\end{verbatim}
\end{small}
\end{minipage}
\caption{String de busca de referência, com dois blocos conectados por \texttt{AND}: o primeiro reúne termos de IA Generativa/LLMs (Intervenção); o segundo, termos de dados governamentais abertos, governo digital e transparência (População, Desfecho, Contexto). A sintaxe é adaptada a cada base no momento da execução.}
\label{fig:string}
\end{figure*}

Antes da execução definitiva, a \textit{string} passou por calibração-piloto. Ela foi confrontada contra um acervo de controle previamente mapeado pela literatura: caso algum desses artigos essenciais não fosse detectado, o escopo de palavras-chave sofreria expansão imediata. % CONFIRMAR: liste 3-5 artigos de controle usados na calibração.

O intervalo temporal abrange de janeiro de 2019 até a data da pesquisa. % CONFIRMAR ano inicial.
A justificativa para o corte repousa no fato de a consolidação dos grandes modelos pré-treinados ter ocorrido sistematicamente a partir de 2019. O esforço apoiou-se exclusivamente no retorno das bases de dados; técnicas auxiliares de prospecção referencial em cascata (\textit{snowballing} para frente e para trás) \citep{wohlin2014} não foram acionadas neste estágio, configurando uma limitação metodológica discutida com maior profundidade na Seção~\ref{sec:ameacas}.

\subsection{Critérios de Seleção e Processo de Triagem}
\label{subsec:selecao}

O filtro do \textit{corpus} bibliográfico baseou-se nos parâmetros expostos na Tabela~\ref{tab:criterios}. O protocolo é excludente: a aprovação final de um artigo exige a conformidade simultânea a todos os critérios de inclusão e a imunidade absoluta a qualquer um dos critérios de exclusão.

\begin{table}[ht]
\centering
\caption{Critérios de inclusão e exclusão}
\label{tab:criterios}
\small
\begin{tabular}{|c|p{6.4cm}|}
\hline
\textbf{ID} & \textbf{Critério} \\ \hline
CI1 & Aplica ou investiga IA Generativa/LLMs em contexto de dados abertos, transparência pública ou governo digital \\ \hline
CI2 & Estudo primário revisado por pares (conferência ou periódico) \\ \hline
CI3 & Publicado na janela temporal definida \\ \hline
CI4 & Texto em inglês ou português \\ \hline
CI5 & Texto completo acessível \\ \hline
CE1 & Aborda apenas um dos eixos (somente IA \textit{ou} somente governo aberto), sem a interseção \\ \hline
CE2 & Estudo secundário (outra revisão, mapeamento ou \textit{survey}) -- considerado apenas como referência para trabalhos relacionados \\ \hline
CE3 & Literatura cinza, editorial, resumo estendido, pôster, \textit{keynote} ou tutorial \\ \hline
CE4 & Duplicata \\ \hline
CE5 & Versão preliminar de estudo posteriormente estendido (mantém-se a versão mais completa) \\ \hline
\end{tabular}
\end{table}

O funil de seleção operou em duas etapas. A triagem inicial ocorreu por meio da leitura rápida de títulos e resumos. Em seguida, os trabalhos pré-selecionados avançaram para a avaliação na íntegra. Todo o referencial foi exportado para o Zotero (responsável pela deduplicação primária) e posteriormente importado para a plataforma Rayyan, onde ocorreu a triagem manual.

Uma única pesquisadora conduziu a revisão. Para mitigar o viés interpretativo individual, todos os parâmetros de inclusão e exclusão foram engessados formalmente no projeto de pesquisa. Além disso, as exclusões no Rayyan exigiram a marcação explícita do rótulo motivador correspondente ao critério CE, gerando uma trilha de auditoria rastreável. O diagrama PRISMA, ilustrado na seção de resultados, consolida todo o fluxo das exclusões.

\subsection{Avaliação de Qualidade}
\label{subsec:qualidade}

Os manuscritos sobreviventes à leitura integral foram então submetidos a um crivo de qualidade focado no rigor científico das evidências. Cada métrica da Tabela~\ref{tab:qualidade} comporta três valores possíveis: $0$ (não atende), $0{,}5$ (atende parcialmente) e $1$ (atende plenamente), culminando em uma pontuação máxima de 5.

Estabeleceu-se o limiar de $2{,}5$ como nota de corte analítica. Trabalhos que igualam ou superam essa barreira compõem o núcleo da síntese informacional sem ressalvas. Aqueles com desempenho inferior não sofrem descarte sumário, mas são apartados em um bloco interpretativo paralelo, exigindo ceticismo deliberado na fase de Discussão.

Esse recurso avaliativo opera também como mecanismo compensatório para publicações em veículos de menor tradição editorial (como simpósios de iniciação ou anais de triagem branda). Pelo critério CI2, estudos primários em eventos oficiais não são descartados na fase inicial. É a matriz de qualidade (notadamente os fatores AQ4 e AQ5) que penaliza a fragilidade argumentativa dessas produções, desidratando sua influência na conclusão final sem violar as regras da filtragem sistemática.

\begin{table}[ht]
\centering
\caption{Critérios de avaliação de qualidade}
\label{tab:qualidade}
\small
\begin{tabular}{|c|p{6.4cm}|}
\hline
\textbf{ID} & \textbf{Critério} \\ \hline
AQ1 & Os objetivos do estudo estão claramente declarados? \\ \hline
AQ2 & O contexto (domínio governamental e dados) está adequadamente descrito? \\ \hline
AQ3 & A abordagem de IA Generativa/LLM está descrita com detalhe suficiente para ser compreendida e replicada? \\ \hline
AQ4 & Os resultados e benefícios são reportados e sustentados por evidências? \\ \hline
AQ5 & As limitações e ameaças à validade são discutidas? \\ \hline
\end{tabular}
\end{table}

\subsection{Extração de Dados}
\label{subsec:extracao}

A captação de evidências seguiu um formulário rígido cujas entradas convergem diretamente para as cinco questões de pesquisa (Tabela~\ref{tab:extracao}). Como complemento à estruturação técnica, levantou-se o inventário bibliométrico padrão (autoria, ano, país e tipologia do veículo), bem como as matrizes metodológicas de validação presentes em cada artigo empírico. Todo o registro das abstrações foi conduzido pela revisora primária, mantendo a indexação rigorosa da informação recolhida com o arquivo de origem para evitar distorções de interpretação fora de contexto.

\begin{table*}[t]
\centering
\caption{Formulário de extração e vínculo com as questões de pesquisa}
\label{tab:extracao}
\small
\renewcommand{\arraystretch}{1.15}
\begin{tabular}{|p{3.6cm}|@{\hspace{0.6em}}p{6.0cm}@{\hspace{0.6em}}|c|}
\hline
\textbf{Campo} & \textbf{Descrição} & \textbf{RQ} \\ \hline
Aplicação / caso de uso & Como a IA Generativa é empregada & RQ1 \\ \hline
Dados governamentais & Tipo e domínio dos dados analisados & RQ2 \\ \hline
Modelo(s) & LLM/IA Generativa, porte, licença & RQ3 \\ \hline
Benefícios & Ganhos de acesso/transparência relatados & RQ4 \\ \hline
Desafios & Limitações técnicas, éticas e de adoção & RQ5 \\ \hline
Avaliação & Método de validação do estudo & --- \\ \hline
\end{tabular}
\end{table*}

\subsection{Síntese dos Dados}
\label{subsec:sintese}

O fechamento da revisão opera sob uma matriz dual. O primeiro pilar baseia-se na \textit{síntese quantitativa descritiva}, responsável por traduzir o volume documental em distribuições macro (como a linha do tempo de publicações, os LLMs mais citados em RQ3 e o espectro de dados analisados em RQ2). O segundo pilar é a \textit{síntese temática}, uma abordagem estritamente qualitativa. Nela, o esforço debruça-se sobre a codificação do texto bruto, extraindo e categorizando de forma orgânica as rotinas governamentais criadas (RQ1), os impactos sociais alcançados (RQ4) e as pedras de tropeço da adoção (RQ5). É o atrito analítico provocado pelo cruzamento dessas duas sínteses que permitirá, no momento da Discussão, organizar o caos de aplicações cívicas de Inteligência Artificial em uma taxonomia coesa e útil para a comunidade científica.

%
%  CLASSE: sbc2025 (apalike-sol). Sem comandos ABNT (\quadro, \citeonline).
%  PACOTES: usa apenas tabular/\hline padrão — não exige booktabs.
%
%  PARÂMETROS A CONFIRMAR (procure por "CONFIRMAR" abaixo):
%   - Janela temporal (ano inicial)
%   - String final após calibração-piloto
%   - Limiares de qualidade
%
%  REFERÊNCIAS .bib:
%   NOVAS (adicionadas a referencias.bib): kitchenham2007, petersen2015,
%         page2021, wohlin2014.
%   NOTA: Kitchenham2004 (mapeamento) já existe no .bib, mas é obra distinta;
%         o Método cita Kitchenham & Charters (2007) = kitchenham2007.
%
%  OBS.: se a classe for de duas colunas e alguma tabela estourar,
%        troque \begin{table} por \begin{table*}.
% =====================================================================

\section{Método}
\label{sec:metodo}

Este trabalho caracteriza-se como uma Revisão Sistemática da Literatura (RSL), método secundário cujo propósito é identificar, avaliar e sintetizar as evidências científicas disponíveis sobre um dado fenômeno de forma reprozível. O protocolo foi planejado e conduzido segundo as diretrizes de Kitchenham e Charters \citeyearpar{kitchenham2007}, complementadas pelas orientações para mapeamento sistemático de Petersen et al. \citeyearpar{petersen2015}. O relato da pesquisa obedece às recomendações do framework PRISMA 2020 \citep{page2021}.

Como as questões levantadas são de caracterização, e não explicativas, optou-se por um desenho \textbf{híbrido}. Isso significa que a busca e a seleção de literatura seguem o rigor procedimental de uma RSL, mas a síntese dos resultados assume a abrangência de um mapeamento de campo. A execução organiza-se em três etapas: planejamento (definição do protocolo e critérios), condução (busca, seleção e extração) e relato (síntese e discussão). Visando proteger a reprodutibilidade do estudo e refrear o viés do pesquisador, todo o protocolo foi consolidado \textit{a priori}.

\subsection{Objetivo e Questões de Pesquisa}
\label{subsec:objetivo_qp}

O objetivo desta revisão é identificar como a Inteligência Artificial Generativa e os Modelos de Linguagem de Grande Porte (\textit{Large Language Models} -- LLMs) vêm sendo empregados para destravar o acesso a dados governamentais abertos e impulsionar a transparência pública. O esforço ancora-se na área de Sistemas de Informação, operando exatamente na interseção entre governo digital, política de dados abertos (\textit{Open Government Data} -- OGD) e apoio à decisão. A partir dessa diretriz central, derivam-se as cinco questões de pesquisa detalhadas na Tabela~\ref{tab:rqs}.

\begin{table*}[t]
\centering
\caption{Questões de pesquisa e foco de extração}
\label{tab:rqs}
\small
\begin{tabular}{|c|p{4.6cm}|p{4.0cm}|}
\hline
\textbf{ID} & \textbf{Questão de pesquisa} & \textbf{Foco da extração} \\ \hline
RQ1 & Quais aplicações de IA Generativa são utilizadas no acesso a dados governamentais abertos e na transparência pública? & Tipo de aplicação / caso de uso \\ \hline
RQ2 & Quais tipos de dados governamentais são analisados? & Natureza e domínio dos dados \\ \hline
RQ3 & Quais LLMs são utilizados? & Modelo(s), porte, aberto/proprietário \\ \hline
RQ4 & Quais benefícios são reportados? & Ganhos de acesso e transparência \\ \hline
RQ5 & Quais desafios e limitações são reportados? & Barreiras técnicas, éticas e de adoção \\ \hline
\end{tabular}
\end{table*}

\subsection{Estratégia de Busca}
\label{subsec:busca}

A varredura automática ocorreu em quatro bases de amplo reconhecimento no domínio da Computação: Scopus, Web of Science, ACM Digital Library e IEEE Xplore. Essa seleção equilibra indexadores multidisciplinares (Scopus, WoS) e repositórios de nicho tecnológico (ACM, IEEE), mitigando o risco de omissão de publicações chave.

Os termos da consulta foram arquitetados a partir das cinco dimensões da estrutura PICOC (Tabela~\ref{tab:picoc}), garantindo que a \textit{string} não recaísse isoladamente sobre a Intervenção tecnológica. O eixo de Intervenção mapeia o maquinário propriamente dito: IA Generativa, modelos de fundação, \textit{transformers}, arquiteturas RAG, agentes conversacionais e LLMs comerciais ou abertos. Em contrapartida, os blocos de População, Desfecho e Contexto amarram o escopo da ferramenta ao setor público, filtrando estritamente estudos sobre portais governamentais, dados legislativos e políticas de transparência digital.

A \textit{string} de referência consta na Figura~\ref{fig:string} e foi adaptada às regras sintáticas de cada base no momento do disparo. Para assegurar aderência ao critério de idioma (CI4), aplicou-se uma variação do comando em português à Biblioteca Digital da SBC/SOL e demais plataformas lusófonas. Com o intuito de distinguir a real escassez de produção acadêmica de meros erros de digitação na busca, a estrutura sintática exata, a data e o volume retornado por base foram documentados no repositório do projeto e recuperados na Seção~\ref{sec:resultados}.

\begin{table*}[t]
\centering
\caption{Termos da busca segundo a estrutura PICOC}
\label{tab:picoc}
\small
\begin{tabular}{|p{2.4cm}|p{6.0cm}|}
\hline
\textbf{Elemento} & \textbf{Termos} \\ \hline
População & Setor público; portais e dados governamentais \\ \hline
Intervenção & IA Generativa; LLMs; modelos de linguagem \\ \hline
Comparação & Não aplicável (revisão de caracterização) \\ \hline
Desfecho & Acesso à informação; transparência; benefícios e desafios \\ \hline
Contexto & Governo digital; \textit{e-government}; OGD \\ \hline
\end{tabular}
\end{table*}

\begin{figure*}[t]
\centering
\begin{minipage}{0.95\textwidth}
\begin{small}
\begin{verbatim}
("generative artificial intelligence" OR "generative AI"
 OR "large language model*" OR "LLM*" OR "foundation model*"
 OR transformer* OR GPT* OR ChatGPT OR Llama OR Mistral OR Gemini
 OR "retrieval-augmented generation" OR RAG
 OR chatbot* OR "conversational agent*")
AND
("open government data" OR "open data" OR "government data"
 OR "public sector data" OR "legislative data" OR "parliamentary data"
 OR "freedom of information" OR FOIA OR "access to information"
 OR transparency OR accountability
 OR "digital government" OR "e-government" OR "open government"
 OR "data portal*")
\end{verbatim}
\end{small}
\end{minipage}
\caption{String de busca de referência, com dois blocos conectados por \texttt{AND}: o primeiro reúne termos de IA Generativa/LLMs (Intervenção); o segundo, termos de dados governamentais abertos, governo digital e transparência (População, Desfecho, Contexto). A sintaxe é adaptada a cada base no momento da execução.}
\label{fig:string}
\end{figure*}

Antes da execução definitiva, a \textit{string} passou por calibração-piloto. Ela foi confrontada contra um acervo de controle previamente mapeado pela literatura: caso algum desses artigos essenciais não fosse detectado, o escopo de palavras-chave sofreria expansão imediata. % CONFIRMAR: liste 3-5 artigos de controle usados na calibração.

O intervalo temporal abrange de janeiro de 2019 até a data da pesquisa. % CONFIRMAR ano inicial.
A justificativa para o corte repousa no fato de a consolidação dos grandes modelos pré-treinados ter ocorrido sistematicamente a partir de 2019. O esforço apoiou-se exclusivamente no retorno das bases de dados; técnicas auxiliares de prospecção referencial em cascata (\textit{snowballing} para frente e para trás) \citep{wohlin2014} não foram acionadas neste estágio, configurando uma limitação metodológica discutida com maior profundidade na Seção~\ref{sec:ameacas}.

\subsection{Critérios de Seleção e Processo de Triagem}
\label{subsec:selecao}

O filtro do \textit{corpus} bibliográfico baseou-se nos parâmetros expostos na Tabela~\ref{tab:criterios}. O protocolo é excludente: a aprovação final de um artigo exige a conformidade simultânea a todos os critérios de inclusão e a imunidade absoluta a qualquer um dos critérios de exclusão.

\begin{table}[ht]
\centering
\caption{Critérios de inclusão e exclusão}
\label{tab:criterios}
\small
\begin{tabular}{|c|p{6.4cm}|}
\hline
\textbf{ID} & \textbf{Critério} \\ \hline
CI1 & Aplica ou investiga IA Generativa/LLMs em contexto de dados abertos, transparência pública ou governo digital \\ \hline
CI2 & Estudo primário revisado por pares (conferência ou periódico) \\ \hline
CI3 & Publicado na janela temporal definida \\ \hline
CI4 & Texto em inglês ou português \\ \hline
CI5 & Texto completo acessível \\ \hline
CE1 & Aborda apenas um dos eixos (somente IA \textit{ou} somente governo aberto), sem a interseção \\ \hline
CE2 & Estudo secundário (outra revisão, mapeamento ou \textit{survey}) -- considerado apenas como referência para trabalhos relacionados \\ \hline
CE3 & Literatura cinza, editorial, resumo estendido, pôster, \textit{keynote} ou tutorial \\ \hline
CE4 & Duplicata \\ \hline
CE5 & Versão preliminar de estudo posteriormente estendido (mantém-se a versão mais completa) \\ \hline
\end{tabular}
\end{table}

O funil de seleção operou em duas etapas. A triagem inicial ocorreu por meio da leitura rápida de títulos e resumos. Em seguida, os trabalhos pré-selecionados avançaram para a avaliação na íntegra. Todo o referencial foi exportado para o Zotero (responsável pela deduplicação primária) e posteriormente importado para a plataforma Rayyan, onde ocorreu a triagem manual.

Uma única pesquisadora conduziu a revisão. Para mitigar o viés interpretativo individual, todos os parâmetros de inclusão e exclusão foram engessados formalmente no projeto de pesquisa. Além disso, as exclusões no Rayyan exigiram a marcação explícita do rótulo motivador correspondente ao critério CE, gerando uma trilha de auditoria rastreável. O diagrama PRISMA, ilustrado na seção de resultados, consolida todo o fluxo das exclusões.

\subsection{Avaliação de Qualidade}
\label{subsec:qualidade}

Os manuscritos sobreviventes à leitura integral foram então submetidos a um crivo de qualidade focado no rigor científico das evidências. Cada métrica da Tabela~\ref{tab:qualidade} comporta três valores possíveis: $0$ (não atende), $0{,}5$ (atende parcialmente) e $1$ (atende plenamente), culminando em uma pontuação máxima de 5.

Estabeleceu-se o limiar de $2{,}5$ como nota de corte analítica. Trabalhos que igualam ou superam essa barreira compõem o núcleo da síntese informacional sem ressalvas. Aqueles com desempenho inferior não sofrem descarte sumário, mas são apartados em um bloco interpretativo paralelo, exigindo ceticismo deliberado na fase de Discussão.

Esse recurso avaliativo opera também como mecanismo compensatório para publicações em veículos de menor tradição editorial (como simpósios de iniciação ou anais de triagem branda). Pelo critério CI2, estudos primários em eventos oficiais não são descartados na fase inicial. É a matriz de qualidade (notadamente os fatores AQ4 e AQ5) que penaliza a fragilidade argumentativa dessas produções, desidratando sua influência na conclusão final sem violar as regras da filtragem sistemática.

\begin{table}[ht]
\centering
\caption{Critérios de avaliação de qualidade}
\label{tab:qualidade}
\small
\begin{tabular}{|c|p{6.4cm}|}
\hline
\textbf{ID} & \textbf{Critério} \\ \hline
AQ1 & Os objetivos do estudo estão claramente declarados? \\ \hline
AQ2 & O contexto (domínio governamental e dados) está adequadamente descrito? \\ \hline
AQ3 & A abordagem de IA Generativa/LLM está descrita com detalhe suficiente para ser compreendida e replicada? \\ \hline
AQ4 & Os resultados e benefícios são reportados e sustentados por evidências? \\ \hline
AQ5 & As limitações e ameaças à validade são discutidas? \\ \hline
\end{tabular}
\end{table}

\subsection{Extração de Dados}
\label{subsec:extracao}

A captação de evidências seguiu um formulário rígido cujas entradas convergem diretamente para as cinco questões de pesquisa (Tabela~\ref{tab:extracao}). Como complemento à estruturação técnica, levantou-se o inventário bibliométrico padrão (autoria, ano, país e tipologia do veículo), bem como as matrizes metodológicas de validação presentes em cada artigo empírico. Todo o registro das abstrações foi conduzido pela revisora primária, mantendo a indexação rigorosa da informação recolhida com o arquivo de origem para evitar distorções de interpretação fora de contexto.

\begin{table*}[t]
\centering
\caption{Formulário de extração e vínculo com as questões de pesquisa}
\label{tab:extracao}
\small
\renewcommand{\arraystretch}{1.15}
\begin{tabular}{|p{3.6cm}|@{\hspace{0.6em}}p{6.0cm}@{\hspace{0.6em}}|c|}
\hline
\textbf{Campo} & \textbf{Descrição} & \textbf{RQ} \\ \hline
Aplicação / caso de uso & Como a IA Generativa é empregada & RQ1 \\ \hline
Dados governamentais & Tipo e domínio dos dados analisados & RQ2 \\ \hline
Modelo(s) & LLM/IA Generativa, porte, licença & RQ3 \\ \hline
Benefícios & Ganhos de acesso/transparência relatados & RQ4 \\ \hline
Desafios & Limitações técnicas, éticas e de adoção & RQ5 \\ \hline
Avaliação & Método de validação do estudo & --- \\ \hline
\end{tabular}
\end{table*}

\subsection{Síntese dos Dados}
\label{subsec:sintese}

O fechamento da revisão opera sob uma matriz dual. O primeiro pilar baseia-se na \textit{síntese quantitativa descritiva}, responsável por traduzir o volume documental em distribuições macro (como a linha do tempo de publicações, os LLMs mais citados em RQ3 e o espectro de dados analisados em RQ2). O segundo pilar é a \textit{síntese temática}, uma abordagem estritamente qualitativa. Nela, o esforço debruça-se sobre a codificação do texto bruto, extraindo e categorizando de forma orgânica as rotinas governamentais criadas (RQ1), os impactos sociais alcançados (RQ4) e as pedras de tropeço da adoção (RQ5). É o atrito analítico provocado pelo cruzamento dessas duas sínteses que permitirá, no momento da Discussão, organizar o caos de aplicações cívicas de Inteligência Artificial em uma taxonomia coesa e útil para a comunidade científica.
